% =============================================================================
\chapter{Wyniki}
\label{ch:results}
% =============================================================================

Ten rozdział przedstawia wyniki przebiegu \code{haiku-v1}. Najpierw opisujemy konfigurację, a potem
kolejno odpowiadamy na pytania badawcze: poprawność (\RQ{1}, \cref{sec:results-correctness}),
rodzaje błędów i defekty benchmarku (\cref{sec:failure-analysis}), samonaprawa (\RQ{2},
\cref{sec:results-repair}), idiomatyczność (\RQ{3}, \cref{sec:results-idiomaticity}) oraz koszty
i wykonywanie instrukcji (\RQ{4}, \cref{sec:cost}).

\section{Konfiguracja eksperymentu}
\label{sec:setup}

\Cref{tab:setup} podsumowuje konfigurację. Każdą z 70 par język--zadanie uruchomiono raz (jedna
próbka). Zadania 1--6 uruchomiono 25~października~2026 na commicie \code{072ea49}. Zadania 7--10
dodano do tego samego przebiegu 13~listopada~2026, po ich scaleniu (commit \code{dc7f1d4}). Ochrona
konfiguracji (\cref{sec:reproducibility}) zapewniła identyczne ustawienia generatora. Wszystkie 70
rozwiązań końcowych oceniono z tą samą konfiguracją sędziego (schemat z uzasadnieniem przed
oceną).

\begin{table}[htbp]
  \centering
  \caption{Konfiguracja przebiegu \code{haiku-v1}.}
  \label{tab:setup}
  \small
  \begin{tabular}{@{}lll@{}}
    \toprule
    & Ustawienie & Wartość \\
    \midrule
    \multirow{5}{*}{Generator}
      & Model & \haiku{} (\code{claude-haiku-5-5}) \\
      & Rozumowanie & adaptacyjne myślenie, wysiłek \code{medium}, podsumowania \\
      & Maks.\ tokenów wyjściowych & 32\,000 na odpowiedź (łącznie z rozumowaniem) \\
      & Rundy naprawy & do 3 (co najwyżej 4 próby na zadanie) \\
      & Próbki na zadanie & 1 \\
    \midrule
    \multirow{4}{*}{Sędzia}
      & Model & \sonnet{} (\code{claude-sonnet-5-5}) \\
      & Rozumowanie & adaptacyjne myślenie, wysiłek \code{high} \\
      & Oceniana próba & ostatnia próba każdego zadania (jeśli się skompilowała) \\
      & Odpowiedź & schemat JSON, uzasadnienie przed ocenami (\cref{lst:schema}) \\
    \midrule
    \multirow{3}{*}{Ewaluacja}
      & Limit czasu & 300\,s na przebieg testów, 4 zadania równolegle \\
      & Skala & 7 języków $\times$ 10 zadań $=$ 70 zadań, 1\,400 testów \\
      & Wywołania API & 98 wywołań generujących, 70 wywołań sędziego \\
    \bottomrule
  \end{tabular}
\end{table}

% -----------------------------------------------------------------------------
\section{RQ1: Poprawność funkcjonalna}
\label{sec:results-correctness}

\Cref{tab:matrix} pokazuje wynik każdego zadania, a \cref{tab:correctness} agreguje go według
języka i paradygmatu. Łącznie model rozwiązał 56 z 70 zadań (80,0\,\%) pierwszą odpowiedzią
i 65 (92,9\,\%) po samonaprawie. W ujęciu testów zaliczył 1\,229 z 1\,400 testów (87,8\,\%) za
pierwszym podejściem i 1\,375 (98,2\,\%) na końcu. W odróżnieniu od pilotażu z HumanEval
(\cref{sec:humaneval}) benchmark pozostawia więc miejsce na obserwację różnic.

\begin{table}[htbp]
  \centering
  \caption{Wyniki dla poszczególnych języków i zadań. Komórka pokazuje liczbę zaliczonych testów
    (z 20) w pierwszej próbie oraz, jeśli jest inna, w próbie ostatniej. \cmark{} oznacza zadanie
    rozwiązane, a indeks górny rundę naprawy, w której je rozwiązano; zacienione komórki z \xmark{}
    nie zostały rozwiązane po trzech rundach naprawy.}
  \label{tab:matrix}
  \footnotesize
  \setlength{\tabcolsep}{3.2pt}
  \begin{tabular}{@{}l*{10}{c}c@{}}
    \toprule
    Język & Z1 & Z2 & Z3 & Z4 & Z5 & Z6 & Z7 & Z8 & Z9 & Z10 & Rozw. \\
    \midrule
    \tabinput{tables/matrix}
    \bottomrule
  \end{tabular}
\end{table}

\begin{table}[htbp]
  \centering
  \caption{Poprawność według języka i paradygmatu: odsetek zaliczonych testów i liczba
    rozwiązanych zadań w pierwszej próbie i po samonaprawie (końcowo).}
  \label{tab:correctness}
  \small
  \begin{tabular}{@{}llrrrr@{}}
    \toprule
    & & \multicolumn{2}{c}{Zaliczone testy (\%)} & \multicolumn{2}{c}{Rozwiązane zadania} \\
    \cmidrule(lr){3-4}\cmidrule(lr){5-6}
    Język & Paradygmat & pierwsza & końcowo & pierwsza & końcowo \\
    \midrule
    C         & imperatywny & 100,0 & 100,0 & 10/10 & 10/10 \\
    Pascal    & imperatywny &  90,0 &  98,0 &  9/10 &  9/10 \\
    Java      & obiektowy   &  70,5 & 100,0 &  6/10 & 10/10 \\
    Smalltalk & obiektowy   &  85,5 &  96,0 &  8/10 &  9/10 \\
    Haskell   & funkcyjny   &  78,0 &  99,5 &  7/10 &  9/10 \\
    Scala     & funkcyjny   &  99,0 & 100,0 &  9/10 & 10/10 \\
    Prolog    & logiczny    &  91,5 &  94,0 &  7/10 &  8/10 \\
    \midrule
    \multicolumn{2}{@{}l}{Imperatywny (C, Pascal)}        & 95,0 & 99,0 & 19/20 & 19/20 \\
    \multicolumn{2}{@{}l}{Obiektowy (Java, Smalltalk)}    & 78,0 & 98,0 & 14/20 & 19/20 \\
    \multicolumn{2}{@{}l}{Funkcyjny (Haskell, Scala)}     & 88,5 & 99,8 & 16/20 & 19/20 \\
    \multicolumn{2}{@{}l}{Logiczny (Prolog)}              & 91,5 & 94,0 &  7/10 &  8/10 \\
    \midrule
    \multicolumn{2}{@{}l}{Łącznie}                        & 87,8 & 98,2 & 56/70 & 65/70 \\
    \bottomrule
  \end{tabular}
\end{table}

\begin{figure}[htbp]
  \centering
  % Share of tests passed at the first and at the final attempt, per language (dumbbell plot).
% Data: results/haiku-v1/summary.csv (first_n_pass, final_n_pass; 200 tests per language).
\begin{tikzpicture}
\begin{axis}[
  reportchart,
  width=0.86\textwidth, height=6.4cm,
  xmin=65, xmax=102, ymin=0.4, ymax=7.6,
  xtick={70,75,80,85,90,95,100},
  xlabel={Zaliczone testy (\%)},
  ytick={1,...,7},
  yticklabels={Prolog (log.), Scala (funk.), Haskell (funk.), Smalltalk (ob.), Java (ob.), Pascal (imp.), C (imp.)},
  xmajorgrids=true, ymajorgrids=false,
  axis x line*=bottom, axis y line*=left,
  y axis line style={draw=none}, ytick style={draw=none},
  legend style={at={(0.0,1.02)}, anchor=south west, legend columns=2, /tikz/every even column/.append style={column sep=10pt}},
  clip=false,
]
  % connectors (first -> final)
  \foreach \y/\a/\b in {7/100/100, 6/90.0/98.0, 5/70.5/100, 4/85.5/96.0, 3/78.0/99.5, 2/99.0/100, 1/91.5/94.0}{
    \edef\temp{\noexpand\draw[draw=neutral, line width=2pt] (axis cs:\a,\y) -- (axis cs:\b,\y);}
    \temp
  }
  \addplot[only marks, mark=*, mark size=3.4pt, mark options={fill=seq250, draw=white, line width=0.8pt}]
    coordinates {(100,7) (90.0,6) (70.5,5) (85.5,4) (78.0,3) (99.0,2) (91.5,1)};
  \addlegendentry{pierwsza próba}
  \addplot[only marks, mark=*, mark size=3.4pt, mark options={fill=seq600, draw=white, line width=0.8pt}]
    coordinates {(100,7) (98.0,6) (100,5) (96.0,4) (99.5,3) (100,2) (94.0,1)};
  \addlegendentry{po samonaprawie (końcowo)}
  % selective direct labels: the two largest repair gains
  \node[font=\scriptsize, text=inksec, anchor=south] at (axis cs:70.5,5.18) {70,5};
  \node[font=\scriptsize, text=inksec, anchor=south] at (axis cs:78.0,3.18) {78,0};
\end{axis}
\end{tikzpicture}

  \caption{Odsetek z 200 testów na język zaliczonych w pierwszej próbie (jasne punkty) i po
    samonaprawie (ciemne). Dla C oba punkty pokrywają się na poziomie 100\,\%.}
  \label{fig:passrate}
\end{figure}

\paragraph{Języki} C jest jedynym językiem, w którym każde zadanie przeszło za pierwszym podejściem,
łącznie ze wszystkimi kontrolami Valgrinda. Dotyczy to również wyniku bez 17 testów przechodzących
na szkielecie (183/183). Następna jest Scala z 9 zadaniami i 99\,\% testów. Java ma najsłabszą
pierwszą próbę: 6 zadań i 70,5\,\% testów (\cref{fig:passrate}). Jej błędy są jednak przeważnie
tanie do naprawienia: dwa z nich to błędy kompilacji spowodowane wymyśloną przez model deklaracją
\code{package} (\cref{sec:failure-analysis}), a wszystkie cztery nierozwiązane za pierwszym razem
zadania w Javie zostały naprawione. Prolog kończy z najniższym wynikiem: 8 rozwiązanych zadań
i 94,0\,\% testów, ale oba jego nierozwiązane zadania są dotknięte defektami w zestawach testów dla
Prologu (\cref{sec:benchmark-defects}).

\paragraph{Paradygmaty} Za pierwszym podejściem kolejność paradygmatów to: imperatywny (19/20 zadań),
funkcyjny (16/20), obiektowy (14/20) i logiczny (7/10). Po samonaprawie paradygmaty imperatywny,
obiektowy i funkcyjny osiągają po 19/20, a logiczny 8/10. Różnice \emph{wewnątrz} paradygmatu są
równie duże jak różnice między paradygmatami. Java rozwiązała za pierwszym podejściem dwa zadania
mniej niż Smalltalk, a Haskell dwa mniej niż Scala. Mniej popularny język z pary nie jest też
konsekwentnie gorszy: Smalltalk wypadł lepiej niż Java, a Pascal niemal dorównał C.

\paragraph{Ostrożność statystyczna} Przy jednej próbce dziesięciu zadań na język żadna z tych różnic
nie jest istotna statystycznie. Największa różnica w pierwszej próbie między dwoma językami, C
(10/10) wobec Javy (6/10), daje $p = 0,087$ w dwustronnym dokładnym teście Fishera. Paradygmat
imperatywny wobec obiektowego (19/20 wobec 14/20) daje $p = 0,091$, a imperatywny wobec logicznego
(19/20 wobec 7/10) $p = 0,095$. Wyniki należy traktować jako przesłanki uzasadniające powtórzone
próbki, a nie jako ustalone efekty (\cref{sec:threats}).

\paragraph{Zadania} Zadanie~1 (logika przestrzenna) okazało się najtrudniejsze: za pierwszym
podejściem rozwiązały je tylko 3 z 7 języków, głównie z powodu danych zdegenerowanych (pojedyncze
punkty, wielokąty puste i płaskie). Następne było zadanie~9 (B-drzewo, 4 z 7), ponieważ dokładny
kształt drzewa nie zostawia miejsca na warianty algorytmu. Zadania~6 (Sudoku) i~8 (TSP) zostały
rozwiązane za pierwszym podejściem we wszystkich siedmiu językach, z zaliczeniem wszystkich 140
testów. Zadanie~2 (Dijkstra) jest jedynym, którego nie rozwiązały dwa języki.

% -----------------------------------------------------------------------------
\section{Analiza błędów}
\label{sec:failure-analysis}

\Cref{tab:outcomes-lang} klasyfikuje wszystkie 98 prób. Występują tylko trzy z dziewięciu kategorii:
nie było żadnego błędu pamięci, awarii, przekroczenia limitu czasu, uciętej odpowiedzi, odmowy ani
pustej odpowiedzi. Wszystkie 98 odpowiedzi zakończyło się normalnie (powód zatrzymania
\code{end_turn}); najdłuższa wykorzystała 23\,034 z 32\,000 dostępnych tokenów wyjściowych.
Najwięcej niezaliczonych testów przypada na Prolog (9 prób), a najwięcej błędów kompilacji na
Javę (4).

\begin{table}[htbp]
  \centering
  \caption{Kategorie wyników wszystkich prób, łącznie z rundami naprawy. Pominięto kategorie, które
    nie wystąpiły (błąd pamięci, awaria w czasie wykonania, przekroczenie czasu, ucięcie, odmowa,
    pusta odpowiedź).}
  \label{tab:outcomes-lang}
  \small
  \begin{tabular}{@{}lrrrr@{}}
    \toprule
    Język & \code{pass} & \code{test_failures} & \code{compile_error} & Próby \\
    \midrule
    C         & 10 & 0 & 0 & 10 \\
    Pascal    &  9 & 3 & 1 & 13 \\
    Java      & 10 & 2 & 4 & 16 \\
    Smalltalk &  9 & 4 & 2 & 15 \\
    Haskell   &  9 & 4 & 2 & 15 \\
    Scala     & 10 & 1 & 1 & 12 \\
    Prolog    &  8 & 9 & 0 & 17 \\
    \midrule
    Łącznie   & 65 & 23 & 10 & 98 \\
    \bottomrule
  \end{tabular}
\end{table}

\paragraph{Błędy kompilacji} Przeanalizowaliśmy wszystkie dziesięć błędów kompilacji
(\cref{tab:compile-errors}). \emph{Połowa z nich w ogóle nie jest błędami programistycznymi}:
odpowiedź zawierała tekst niebędący kodem. W czterech rundach naprawy (Java, zadania 1 i 9; Scala,
zadanie 1; Smalltalk, zadanie 9) model zaczął odpowiedź od zdania wyjaśniającego błąd
(\cref{lst:leak}), mimo że prompt naprawczy mówi „do NOT write any explanations”. W jednej
pierwszej próbie (Haskell, zadanie 4) kod kończył się zabłąkanym znacznikiem HTML \code{</br>}.
Ponieważ te odpowiedzi nie miały znaczników bloków Markdown, mechanizm wycinania pozostawił je bez
zmian, a kompilator je odrzucił (\cref{sec:extraction}). Pozostałe pięć błędów kompilacji to
rzeczywiste pomyłki. Są to: deklaracja \code{package} niezgodna z domyślnym pakietem kontraktu
(Java, zadania 9 i 10), lokalna zmienna \code{L} pisana wielką literą, którą Haskell interpretuje
jako konstruktor danych (zadanie 3), wywołanie \code{IsInfinity} bez modułu \code{Math} (Pascal,
zadanie 5) oraz rozwiązanie w Smalltalku, którego klasy się nie załadowały (zadanie 2).

\begin{table}[htbp]
  \centering
  \caption{Przyczyny dziesięciu błędów kompilacji.}
  \label{tab:compile-errors}
  \small
  \begin{tabularx}{\textwidth}{@{}llcY@{}}
    \toprule
    Język & Zadanie & Próba & Przyczyna \\
    \midrule
    Java      & 1  & 1 & odpowiedź zaczyna się od wyjaśnienia prozą \\
    Java      & 9  & 1 & odpowiedź zaczyna się od wyjaśnienia prozą \\
    Scala     & 1  & 1 & odpowiedź zaczyna się od wyjaśnienia prozą \\
    Smalltalk & 9  & 1 & odpowiedź zaczyna się od wyjaśnienia prozą \\
    Haskell   & 4  & 0 & zabłąkany znacznik \code{</br>} po ostatniej linii kodu \\
    \midrule
    Java      & 9  & 0 & wymyślona deklaracja \code{package com.example.btree;} ukrywa klasę kontraktu \\
    Java      & 10 & 0 & wymyślona deklaracja \code{package com.example;} ukrywa klasę kontraktu \\
    Haskell   & 3  & 0 & lokalne wiązanie o nazwie \code{L}, interpretowane jako konstruktor danych \\
    Pascal    & 5  & 0 & \code{IsInfinity} użyte bez modułu \code{Math} \\
    Smalltalk & 2  & 0 & klasy nie załadowały się do obrazu (\texttt{nil doesNotUnderstand: \#new}) \\
    \bottomrule
  \end{tabularx}
\end{table}

\begin{lstlisting}[style=prompt,caption={Początek dwóch odpowiedzi naprawczych, które zostały zapisane jako pliki źródłowe (Java, zadanie 9; Smalltalk, zadanie 9; linie skrócone).},label={lst:leak}]
package com.example.btree is the root cause. BTree.java and Tests.java are in the default package, so BTreeImpl in a named package cannot see BTree and its inherited nested type Node. ...

The failure is in the deletion path. Every test that removes a key fails, while insertion and the absent-key test pass. The likely cause is `OrderedCollection>>...
\end{lstlisting}

Jest to ten sam problem, który uczynił modele lokalne bezużytecznymi (\cref{sec:local-llms}), ale
o znacznie mniejszej częstości: 4 z 28 odpowiedzi naprawczych (14\,\%) i 1 z 70 pierwszych
odpowiedzi. Prompt naprawczy wydaje się skłaniać do wyjaśnień, ponieważ model właśnie usłyszał, że
jego kod zawiódł. Te wycieki prozy zawyżają też liczbę błędów kompilacji w Javie, Scali i Smalltalku
i kosztują każdy z tych języków jedną rundę naprawy.

\paragraph{Zadania nierozwiązane} Pięć zadań pozostało nierozwiązanych:
\begin{itemize}
  \item \emph{Pascal, zadanie 5 (16/20).} Po dodaniu brakującego modułu cztery testy nadal nie
    przechodziły: zagnieżdżone operatory jednoargumentowe, jednoargumentowy minus a potęgowanie
    ($-2^2 = -(2^2)$), funkcje wbudowane i błędy dziedziny. Są to reguły priorytetów i semantyki
    błędów ze specyfikacji.
  \item \emph{Smalltalk, zadanie 2 (12/20).} Odtwarzanie ścieżki przechodzi, ale zapytania
    o odległości nie. Program testowy w Smalltalku przechwytuje każdy wyjątek i wypisuje tylko
    \code{[FAIL]}, więc prompty naprawcze nie zawierały żadnej diagnostyki poza nazwami testów.
    Liczba zaliczonych testów wzrosła z 3 do 12, po czym przestała się zmieniać.
  \item \emph{Haskell, zadanie 3 (19/20).} Jeden test, obsługa brakujących argumentów
    (\code{NULL}), nadal nie przechodził po trzech rundach.
  \item \emph{Prolog, zadania 1 i 2 (11/20 i 17/20).} Oba są dotknięte defektami w zestawach testów
    dla Prologu, opisanymi poniżej.
\end{itemize}

\subsection{Defekty w zestawach testów dla Prologu}
\label{sec:benchmark-defects}

Analiza dwóch nierozwiązanych zadań w Prologu wykazała, że błędy są, przynajmniej częściowo,
spowodowane przez benchmark, a nie przez model.

\paragraph{Zadanie 2: błąd o jeden w funkcji pomocniczej testów} Testy obciążeniowe A10, B05 i B10
budują łańcuch o 1\,000 wierzchołków za pomocą predykatu pomocniczego \code{create_linear_chain/3}.
Wersja dla Prologu iteruje po wszystkich wierzchołkach $0,\dots,N{-}1$ i dodaje krawędź od każdego
z nich do następnika, w tym krawędź $999 \to 1000$, która wskazuje poza graf o $N = 1000$
wierzchołkach. Funkcja pomocnicza w C kończy o jeden wierzchołek wcześniej (\cref{lst:defect}).
Specyfikacja dla Prologu nie mówi, co \code{add_edge/5} ma zrobić z wierzchołkiem spoza zakresu.
Implementacja modelu go odrzuca, co jest rozsądnym wyborem, więc graf testowy nigdy nie zostaje
zbudowany i wszystkie trzy testy nie przechodzą niezależnie od implementacji algorytmu Dijkstry.
Prompty naprawcze pokazywały tylko nazwy testów, więc model nie mógł odkryć przyczyny. Liczba
zaliczonych testów pozostała równa 17 we wszystkich czterech próbach. Pozostałe 17 testów przeszło
za pierwszym podejściem; czy rozwiązanie przechodzi też testy obciążeniowe, będzie wiadomo dopiero
po poprawieniu funkcji pomocniczej.

\begin{lstlisting}[language=Prolog,caption={Predykat pomocniczy budujący łańcuch w testach dla Prologu (u góry) dodaje krawędź spoza zakresu $N{-}1 \to N$; funkcja pomocnicza w C (u dołu) tego nie robi.},label={lst:defect}]
create_linear_chain(NumVertices, Weight, Graph) :-
    create_graph(NumVertices, G0),
    Max is NumVertices - 1,
    numlist(0, Max, Indices),                     % 0 .. N-1
    foldl(add_chain_edge(Weight), Indices, G0, Graph).
add_chain_edge(Weight, I, GIn, GOut) :-
    J is I + 1,                                   % for I = N-1: J = N (out of range)
    add_edge(GIn, I, J, Weight, GOut).

/* C: for (int i = 0; i < num_vertices - 1; i++) graph_add_edge(g, i, i + 1, weight); */
\end{lstlisting}

\paragraph{Zadanie 1: niedospecyfikowany kontrakt} W odróżnieniu od późniejszych kontraktów dla
Prologu (\cref{lst:contract-pl}) szkielet zadania 1 nie dokumentuje ani nazw argumentów, ani ich
trybów: \code{check_dynamic_collision(_, _, _)}. Jego komentarz zapisuje też strukturę termów
z nazwami pisanymi wielką literą (\code{Polygon(...)}), które różnią się od termów używanych
w testach (\code{polygon(...)}). Testy wywołują \code{check_dynamic_collision(Obstacle, Robot,
Velocity)}, zgodnie z nagłówkiem w C, podczas gdy specyfikacja w języku naturalnym wymienia robota
jako pierwszego. Model wybrał kolejność \code{(Robot, Obstacle, Velocity)}, która zamienia miejscami
wielokąt ruchomy i nieruchomy, przez co wszystkie osiem nietrywialnych testów kolizji (B01--B08) nie
przechodzi. Test A08 oczekuje ponadto, że argumenty \code{none} dadzą wynik \code{none}, o czym
specyfikacja dla Prologu nie wspomina. W C, Javie, Haskellu i Scali nazwy lub typy parametrów
rozstrzygają kolejność wywołania i żaden inny język nie oblał testów kolizji.

\begin{findingbox}[Spostrzeżenie.]
Oba nierozwiązane zadania w Prologu pochodzą z wczesnych, ręcznie napisanych zestawów testów;
generowane zadania 6--10 nie mają takich problemów. Surowe liczby w tym rozdziale podajemy tak, jak
zostały zmierzone, ale wyniki Prologu dla zadań 1 i 2 zaniżają możliwości modelu. Zestawy te należy
poprawić przed kolejnymi przebiegami (\cref{ch:conclusions}).
\end{findingbox}

% -----------------------------------------------------------------------------
\section{RQ2: Samonaprawa}
\label{sec:results-repair}

\Cref{tab:repair} i \cref{fig:repair} pokazują, kiedy rozwiązano każde zadanie. Spośród 14 zadań
nierozwiązanych za pierwszym podejściem naprawiono 9 (64\,\%): 5 w pierwszej rundzie i 4 w drugiej.
Trzecia runda nigdy nie pomogła. Naprawione zadanie wymagało średnio 1,4 rundy.

\begin{table}[htbp]
  \centering
  \caption{Samonaprawa według języka. „Naprawione” to zadania, które nie przeszły w próbie 0, ale
    zostały rozwiązane później; trajektoria podaje liczbę zaliczonych testów w kolejnych próbach
    (C = błąd kompilacji, P = zaliczone).}
  \label{tab:repair}
  \small
  \begin{tabularx}{\textwidth}{@{}lccccY@{}}
    \toprule
    Język & W próbie 0 & Naprawione & Nierozw. & Śr.\ runda & Trajektorie nieudanych zadań \\
    \midrule
    C         & 10 & --       & 0 & --  & -- \\
    Pascal    &  9 & 0/1      & 1 & --  & Z5: 0C, 16, 16, 16 \\
    Java      &  6 & 4/4      & 0 & 1,5 & Z1: 19, 0C, 20P; Z7: 2, 20P; Z9: 0C, 0C, 20P; Z10: 0C, 20P \\
    Smalltalk &  8 & 1/2      & 1 & 2,0 & Z2: 0C, 3, 12, 12; Z9: 11, 0C, 20P \\
    Haskell   &  7 & 2/3      & 1 & 1,0 & Z1: 16, 20P; Z3: 0C, 19, 19, 19; Z4: 0C, 20P \\
    Scala     &  9 & 1/1      & 0 & 2,0 & Z1: 18, 0C, 20P \\
    Prolog    &  7 & 1/3      & 2 & 1,0 & Z1: 10, 11, 11, 11; Z2: 17, 17, 17, 17; Z9: 16, 20P \\
    \midrule
    Łącznie   & 56 & 9/14     & 5 & 1,4 & \\
    \bottomrule
  \end{tabularx}
\end{table}

\begin{figure}[htbp]
  \centering
  % Attempt at which each task was solved, per language (stacked bars, 10 tasks per language).
% Data: results/haiku-v1/report.md, section 4 (self-repair).
\begin{tikzpicture}
\begin{axis}[
  reportchart,
  xbar stacked,
  width=0.86\textwidth, height=6.2cm,
  bar width=10pt,
  xmin=0, xmax=10, xtick={0,2,...,10},
  xlabel={Zadania (z 10)},
  ytick={1,...,7}, ymin=0.4, ymax=7.6,
  yticklabels={Prolog, Scala, Haskell, Smalltalk, Java, Pascal, C},
  xmajorgrids=true,
  axis x line*=bottom, axis y line*=left,
  y axis line style={draw=none}, ytick style={draw=none},
  legend style={at={(0.0,1.02)}, anchor=south west, legend columns=4,
                /tikz/every even column/.append style={column sep=8pt}},
  every axis plot/.append style={draw=white, line width=0.8pt},
  legend image code/.code={\draw[#1, draw=none] (0cm,-0.1cm) rectangle (0.3cm,0.12cm);},
]
  \addplot[fill=seq600] coordinates {(10,7) (9,6) (6,5) (8,4) (7,3) (9,2) (7,1)};
  \addlegendentry{próba 0 (pierwsza odpowiedź)}
  \addplot[fill=seq400] coordinates {(0,7) (0,6) (2,5) (0,4) (2,3) (0,2) (1,1)};
  \addlegendentry{runda naprawy 1}
  \addplot[fill=seq250] coordinates {(0,7) (0,6) (2,5) (1,4) (0,3) (1,2) (0,1)};
  \addlegendentry{runda naprawy 2}
  \addplot[fill=neutral] coordinates {(0,7) (1,6) (0,5) (1,4) (1,3) (0,2) (2,1)};
  \addlegendentry{nierozwiązane}
\end{axis}
\end{tikzpicture}

  \caption{Próba, w której rozwiązano każde z dziesięciu zadań w danym języku.}
  \label{fig:repair}
\end{figure}

\paragraph{Co udaje się naprawić} Naprawa działa najlepiej, gdy informacja zwrotna wskazuje
bezpośrednio na defekt. Poza Prologiem wszystkie pięć zadań, które w pierwszej próbie oblały testy,
zostało naprawionych w ciągu jednej lub dwóch rund na podstawie linii \code{[FAIL]}
i komunikatów programu testowego; w Prologu udało się to w jednym przypadku na trzy. Z sześciu błędów
kompilacji w pierwszej próbie pięć skompilowało się po jednej rundzie, ponieważ komunikat
kompilatora nazywa problem; szósty (Java, zadanie 9) skompilował się po drugiej rundzie, gdyż
pierwsza odpowiedź naprawcza zaczynała się od prozy. Tylko trzy z sześciu przeszły następnie
wszystkie testy. Pozostałe błędy miały charakter semantyczny (reguły priorytetów w Pascalu) lub
były niewidoczne w informacji zwrotnej.

\paragraph{Kiedy naprawa utyka} We wszystkich pięciu nierozwiązanych zadaniach liczba zaliczonych
testów przestała się zmieniać po co najwyżej dwóch rundach (\cref{tab:repair}). W każdym przypadku
informacja zwrotna niosła niewiele treści. Defektów w Prologu nie da się dostrzec z nazw testów,
program testowy w Smalltalku połyka wyjątki, a jedyny pozostały błąd w Haskellu dotyczy konwencji
(brakujących argumentów), której log nie wyjaśnia. Jest to zgodne z ustaleniami
\citet{olausson2024selfrepair}, którzy wskazali jakość informacji zwrotnej jako wąskie gardło
samonaprawy. Dla twórców benchmarków oznacza to, że programy testowe powinny wypisywać wartość
oczekiwaną i rzeczywistą dla niespełnionej asercji.

\paragraph{Regresje} Trzy naprawy przeszły przez błąd kompilacji, zanim się powiodły (Java,
zadanie 1; Scala, zadanie 1; Smalltalk, zadanie 9). W każdej z nich pierwsza odpowiedź naprawcza
zaczynała się od prozy (\cref{tab:compile-errors}), a druga ją usunęła. Bez tych błędów formatu
zadania te zostałyby prawdopodobnie naprawione o rundę wcześniej, a średnia runda naprawy byłaby
niższa.

\paragraph{Paradygmaty} Odsetki naprawy dla paradygmatów wynoszą 0/1 (imperatywny), 5/6 (obiektowy),
3/4 (funkcyjny) i 1/3 (logiczny). Liczby są zbyt małe, by je porównywać. Niski wynik paradygmatu
logicznego tłumaczą dwa defekty w zestawach testów, a wynik obiektowy jest zdominowany przez tanie
poprawki błędów kompilacji w Javie.

% -----------------------------------------------------------------------------
\section{RQ3: Idiomatyczność}
\label{sec:results-idiomaticity}

\Cref{tab:judge} zestawia średnie oceny sędziego według języka i paradygmatu, a \cref{fig:judge}
porównuje paradygmaty.

\begin{table}[htbp]
  \centering
  \caption{Średnie oceny idiomatyczności rozwiązań końcowych (1 = słabo, 5 = doskonale), wystawione
    na ślepo przez \sonnet{}. $n$ = liczba ocenionych rozwiązań; AW = średnia liczba antywzorców
    wskazanych na rozwiązanie. Najlepsza wartość w kolumnie spośród języków jest pogrubiona.}
  \label{tab:judge}
  \small
  \begin{tabular}{@{}lrccccc r@{}}
    \toprule
    & $n$ & Paradygmat & Język & Czytelność & Projekt & Średnia & AW \\
    \midrule
    C         & 10 & 4,10 & 3,50 & 3,60 & 3,40 & 3,65 & 4,6 \\
    Pascal    & 10 & 4,00 & 3,40 & 3,50 & 3,50 & 3,60 & 4,6 \\
    Java      & 10 & 3,50 & 3,80 & 3,70 & 3,60 & 3,65 & 4,0 \\
    Smalltalk & 10 & 3,00 & 3,30 & 3,50 & 3,20 & 3,25 & 4,8 \\
    Haskell   & 10 & 4,20 & 3,70 & 3,70 & \textbf{3,70} & 3,83 & 3,9 \\
    Scala     & 10 & \textbf{4,30} & \textbf{4,00} & \textbf{3,90} & \textbf{3,70} & \textbf{3,98} & \textbf{3,1} \\
    Prolog    & 10 & 3,40 & 3,40 & 3,50 & 3,30 & 3,40 & 5,0 \\
    \midrule
    Imperatywny & 20 & 4,05 & 3,45 & 3,55 & 3,45 & 3,62 & 4,6 \\
    Obiektowy   & 20 & 3,25 & 3,55 & 3,60 & 3,40 & 3,45 & 4,4 \\
    Funkcyjny   & 20 & 4,25 & 3,85 & 3,80 & 3,70 & 3,90 & 3,5 \\
    Logiczny    & 10 & 3,40 & 3,40 & 3,50 & 3,30 & 3,40 & 5,0 \\
    \bottomrule
  \end{tabular}
\end{table}

\begin{figure}[htbp]
  \centering
  % Mean judge scores per paradigm and criterion (1-5 scale; 20 solutions per paradigm, 10 for logic).
% Data: results/haiku-v1/report.md, section 5. Colours: categorical slots 1-4 in fixed order.
\begin{tikzpicture}
\begin{axis}[
  reportchart,
  ybar=1.5pt,
  width=0.9\textwidth, height=6.2cm,
  bar width=11pt,
  ymin=1, ymax=5, ytick={1,2,3,4,5},
  ylabel={Średnia ocena (1--5)},
  symbolic x coords={P, L, R, D},
  xtick=data,
  xticklabels={Idiomy paradygmatu, Idiomy języka, Czytelność, Projekt},
  enlarge x limits=0.16,
  ymajorgrids=true,
  axis x line*=bottom, axis y line*=left,
  x axis line style={draw=inkmuted},
  legend style={at={(0.0,1.02)}, anchor=south west, legend columns=4,
                /tikz/every even column/.append style={column sep=8pt}},
  every axis plot/.append style={draw=white, line width=0.6pt},
  legend image code/.code={\draw[#1, draw=none] (0cm,-0.1cm) rectangle (0.3cm,0.12cm);},
]
  \addplot[fill=cat1] coordinates {(P,4.05) (L,3.45) (R,3.55) (D,3.45)};
  \addlegendentry{imperatywny}
  \addplot[fill=cat2] coordinates {(P,3.25) (L,3.55) (R,3.60) (D,3.40)};
  \addlegendentry{obiektowy}
  \addplot[fill=cat3] coordinates {(P,4.25) (L,3.85) (R,3.80) (D,3.70)};
  \addlegendentry{funkcyjny}
  \addplot[fill=cat4] coordinates {(P,3.40) (L,3.40) (R,3.50) (D,3.30)};
  \addlegendentry{logiczny}
\end{axis}
\end{tikzpicture}

  \caption{Średnie oceny sędziego według paradygmatu i kryterium. Oś zaczyna się od 1, najniższej
    oceny w rubryce.}
  \label{fig:judge}
\end{figure}

\paragraph{Kod funkcyjny jest najbardziej idiomatyczny} Paradygmat funkcyjny prowadzi we wszystkich
kryteriach (średnia 3,90), przy czym Scala jest najlepszym językiem (3,98) i ma najmniej
antywzorców na rozwiązanie (3,1). Dalej plasuje się paradygmat imperatywny (3,62), następnie
obiektowy (3,45) i logiczny (3,40). Smalltalk ma najniższą ocenę spośród wszystkich języków (3,25).

\paragraph{Idiomatyczność paradygmatu najsilniej różnicuje paradygmaty} Kryterium mierzące sam
paradygmat ma największy rozrzut, od 3,00 (Smalltalk) do 4,30 (Scala). Kod imperatywny (4,05)
i funkcyjny (4,25) uzyskuje wysokie oceny, obiektowy (3,25) i logiczny (3,40) niskie. Smalltalk
otrzymał za paradygmat dokładnie 3 we wszystkich dziesięciu zadaniach (\cref{tab:paradigm-scores}
w \cref{app:results}). Antywzorce wskazane przez sędziego (\cref{tab:antipatterns}) pokazują
dlaczego. W językach obiektowych model pisze kod proceduralny w klasach: statyczne metody pomocnicze
zamiast obiektów dziedzinowych w Javie oraz funkcje geometryczne w klasie usługowej z jawnymi
sprawdzeniami \code{isNil ifTrue:ifFalse:} w Smalltalku. W Prologu emuluje konstrukcję
if-then-else odcięciami, a nawet zgaduje kolejność argumentów za pomocą klauzul chronionych
odcięciem. Mówiąc wprost, model najpierw pisze algorytm, a dopiero potem dopasowuje go do języka.
Sprawdza się to w językach imperatywnych, gdzie algorytm \emph{jest} idiomem, oraz w językach
funkcyjnych, gdzie zwinięcia i dopasowanie wzorców naturalnie odwzorowują te algorytmy. Sprawdza
się gorzej tam, gdzie paradygmat wymaga innej dekompozycji (obiekty, relacje).

\begin{table}[htbp]
  \centering
  \caption{Antywzorce wskazane przez sędziego (po trzy przykłady na język z \file{report.md}).}
  \label{tab:antipatterns}
  \small
  \begin{tabularx}{\textwidth}{@{}lY@{}}
    \toprule
    Język & Antywzorce \\
    \midrule
    C         & ponowna implementacja \code{fabs} zamiast użycia \code{math.h}; jedna długa funkcja łącząca generowanie punktów, usuwanie duplikatów, budowę otoczki i kopiowanie; ręczna pętla kopiująca zamiast \code{memcpy} \\
    Pascal    & własna implementacja quicksorta zamiast sortowania z RTL lub generycznego; zbędne ręczne czyszczenie pól zarządzanych przed \code{Dispose}; alokacja natychmiast zastępowana inną \\
    Java      & statyczne metody pomocnicze zamiast obiektów dziedzinowych (brak typów otoczki/odcinka); zduplikowane prymitywy geometryczne; surowy bufor \code{Point[]} w stylu C z ręczną żonglerką indeksami \\
    Smalltalk & zduplikowana budowa dolnej i górnej otoczki; funkcje geometryczne w klasie usługowej zamiast w obiektach dziedzinowych; \code{isNil ifTrue:ifFalse:} zamiast \code{ifNil:} \\
    Haskell   & częściowa funkcja \code{head} w \code{map head . groupBy}; \code{length xs < 3} zamiast dopasowania wzorca; leniwe \code{foldl} zamiast \code{foldl'} \\
    Scala     & krawędzie budowane arytmetyką indeksów i modulo zamiast złączenia (\code{zip}) z obróconym wektorem; niespójne użycie epsilona; modyfikowalne \code{var}/\code{while} w wyszukiwaniu binarnym zamiast rekursji ogonowej \\
    Prolog    & odcięcie po warunku dla emulacji if-else; \code{append(A0,[_],A)} do usunięcia ostatniego elementu; pobłażliwe zgadywanie kolejności argumentów w \code{create_polygon/3} z użyciem odcięcia \\
    \bottomrule
  \end{tabularx}
\end{table}

\paragraph{Poprawność i idiomatyczność są w dużej mierze niezależne} Java rozwiązała wszystkie
dziesięć zadań, ale za idiomy paradygmatu uzyskała tylko 3,50. C i Scala są w pełni poprawne
i oceniane wysoko, natomiast Smalltalk i Prolog wypadają nisko w obu wymiarach. Rozwiązania, które
przeszły za pierwszym podejściem, uzyskują średnio 3,64, a te, które wymagały naprawy lub pozostały
nierozwiązane, 3,55, co jest różnicą pomijalną. Poprawność i idiomatyczność trzeba zatem mierzyć
osobno.

\paragraph{Idiomatyczność języka podąża za jego popularnością} W kryterium języka prowadzą szeroko
stosowane Scala (4,00), Java (3,80) i Haskell (3,70), a Smalltalk (3,30), Pascal (3,40) i Prolog
(3,40) zamykają stawkę. Jest to spójne z mniejszą ilością danych treningowych dla tych
języków~\cite{cassano2024knowledge}: model słabiej zna ich biblioteki standardowe i konwencje
(\code{ifNil:}, sortowanie z RTL, predykaty biblioteczne).

\paragraph{Wpływ zadania} Ściśle wyspecyfikowane zadania 7--10 uzyskują niższe oceny (3,39--3,50)
niż zadania 1--6 (3,61--3,93; \cref{tab:judge-heatmap} w \cref{app:results}). Narzucają one
dokładne procedury (operacje na B-drzewie według CLRS, reguły remisów w algorytmie Prima, stabilny
podział trójdzielny), które pozostawiają mało miejsca na idiomatyczną przebudowę. Należy o tym
pamiętać przy porównywaniu lub dodawaniu zadań.

\paragraph{Rozdzielczość ocen} Sędzia korzystał z wąskiej części skali. Wszystkie 280 ocen mieści
się między 3 a 5. Ocenę 5 wystawiono tylko siedem razy, zawsze za idiomatyczność paradygmatu (cztery
w Scali, dwie w Haskellu, jedna w C), a żadne kryterium nigdy nie otrzymało 1 ani 2.
Idiomatyczność języka, czytelność i projekt zawsze otrzymywały 3 lub 4. Przy dziesięciu
rozwiązaniach na język zmiana jednej oceny przesuwa średnią kryterium dla języka o 0,1, więc różnic
rzędu 0,1--0,2 nie należy nadinterpretowywać.

\paragraph{Wiarygodność sędziego} W ramach kalibracji oceniono dwukrotnie pięć rozwiązań wzorcowych
i trzy rozwiązania z przebiegu. Wszystkie 32 pary powtórzonych ocen były identyczne. Między
kalibracją, która używała wcześniejszego schematu bez uzasadnień dla poszczególnych kryteriów,
a ostatecznym sędzią z uzasadnieniem przed oceną trzy rozwiązania z przebiegu różniły się w 1 z 12
ocen (czytelność rozwiązania zadania 3 w C, 5 wobec 4). Sędzia jest zatem stabilny. Czy jest
\emph{trafny}, to osobna kwestia. Rozwiązania wzorcowe napisane na potrzeby walidacji zestawów
testów uzyskały średnio 3,05, wyraźnie mniej niż własne rozwiązania modelu \haiku{} dla tych samych
pięciu zadań (3,90). Może to po prostu oznaczać, że wzorce pisano w celu walidacji testów, a nie
jako rozwiązania modelowe. Może to być też preferowanie własnych odpowiedzi: generator i sędzia
pochodzą z tej samej rodziny modeli, a sędziowie LLM, jak wiadomo, faworyzują tekst podobny do
własnego~\cite{panickssery2024selfpreference}. Bez ocen ludzkich nie potrafimy rozstrzygnąć między
tymi wyjaśnieniami (\cref{sec:threats}).

% -----------------------------------------------------------------------------
\section{RQ4: Koszt, wydajność i wykonywanie instrukcji}
\label{sec:cost}

\Cref{tab:efficiency} i \cref{fig:size} podsumowują koszt przebiegu, rozmiar wygenerowanego kodu
i zgodność modelu z wymaganym formatem odpowiedzi.

\begin{table}[htbp]
  \centering
  \caption{Wydajność i wykonywanie instrukcji według języka. LOC = mediana liczby linii kodu
    pierwszej odpowiedzi; tokeny i czas generowania to średnie na wywołanie API (tokeny wyjściowe
    obejmują rozumowanie); czas ewaluacji to średni czas jednego przebiegu testów w Dockerze;
    naruszenia formatu to odpowiedzi ze znacznikami bloków Markdown.}
  \label{tab:efficiency}
  \small
  \begin{tabular}{@{}lrrrrrrr@{}}
    \toprule
    Język & LOC & \makecell[r]{Tokeny\\wyjściowe} & \makecell[r]{Czas\\gen.\ (s)} & \makecell[r]{Czas\\ewal.\ (s)} & Wywołania & \makecell[r]{Koszt\\(USD)} & \makecell[r]{Naruszenia\\formatu} \\
    \midrule
    C         & 248,5 & 6\,014 & 24,1 &  2,2 & 10 & 0,032 &  5/10 \\
    Pascal    & 273,5 & 9\,320 & 35,5 &  1,8 & 13 & 0,069 &  0/13 \\
    Java      & 184,5 & 3\,791 & 14,0 &  2,6 & 16 & 0,037 &  9/16 \\
    Smalltalk & 132,0 & 7\,166 & 29,8 &  1,2 & 15 & 0,062 &  6/15 \\
    Haskell   & 127,0 & 4\,762 & 19,2 & 32,1 & 15 & 0,042 &  1/15 \\
    Scala     & 104,5 & 4\,459 & 17,3 & 29,3 & 12 & 0,030 &  6/12 \\
    Prolog    & 160,0 & 9\,463 & 39,2 & 26,2 & 17 & 0,095 & 10/17 \\
    \midrule
    Łącznie   & --    & 6\,482 & 25,9 & --   & 98 & 0,368 & 37/98 \\
    \bottomrule
  \end{tabular}
\end{table}

\begin{figure}[htbp]
  \centering
  % Two small multiples: code size and output tokens per call, per language.
% Data: results/haiku-v1/attempts.jsonl (code_lines of attempt 0; output_tokens of all calls).
\begin{subfigure}[t]{0.49\textwidth}
\centering
\begin{tikzpicture}
\begin{axis}[
  reportchart, xbar, width=0.93\linewidth, height=5.6cm, bar width=8pt,
  xmin=0, xmax=300, xtick={0,100,200,300},
  xlabel={Mediana liczby linii kodu (próba 0)},
  ytick={1,...,7}, ymin=0.4, ymax=7.6,
  yticklabels={Prolog, Scala, Haskell, Smalltalk, Java, Pascal, C},
  xmajorgrids=true, axis x line*=bottom, axis y line*=left,
  y axis line style={draw=none}, ytick style={draw=none},
]
  \addplot[fill=cat1, draw=none] coordinates {(248.5,7) (273.5,6) (184.5,5) (132,4) (127,3) (104.5,2) (160,1)};
\end{axis}
\end{tikzpicture}
\caption{Rozmiar rozwiązania.}
\label{fig:size-loc}
\end{subfigure}\hfill
\begin{subfigure}[t]{0.49\textwidth}
\centering
\begin{tikzpicture}
\begin{axis}[
  reportchart, xbar, width=0.93\linewidth, height=5.6cm, bar width=8pt,
  xmin=0, xmax=10, xtick={0,2,...,10},
  xlabel={Tokeny wyjściowe na wywołanie, tys.\ (z rozumowaniem)},
  ytick={1,...,7}, ymin=0.4, ymax=7.6,
  yticklabels={Prolog, Scala, Haskell, Smalltalk, Java, Pascal, C},
  xmajorgrids=true, axis x line*=bottom, axis y line*=left,
  y axis line style={draw=none}, ytick style={draw=none},
]
  \addplot[fill=cat1, draw=none] coordinates {(6.014,7) (9.320,6) (3.791,5) (7.166,4) (4.762,3) (4.459,2) (9.463,1)};
\end{axis}
\end{tikzpicture}
\caption{Nakład na generowanie.}
\label{fig:size-tokens}
\end{subfigure}

  \caption{Rozmiar wygenerowanych rozwiązań i nakład na generowanie według języka. Są to dwie
    odrębne miary, stąd dwie odrębne skale.}
  \label{fig:size}
\end{figure}

\paragraph{Zwięzłość} Rozwiązania funkcyjne są najkrótsze. Mediana rozwiązania w Scali ma 105 linii,
a w Haskellu 127, wobec 249 w C i 274 w Pascalu, tj.\ rozwiązania imperatywne są 2--2,6 razy
dłuższe przy tej samej funkcjonalności. Smalltalk (132) i Prolog (160) plasują się pośrodku.

\paragraph{Nakład na generowanie} Liczba tokenów wyjściowych, łącznie z rozumowaniem, jest
najwyższa dla Prologu (9\,463 na wywołanie) i Pascala (9\,320), a najniższa dla Javy (3\,791). Kod
w Pascalu jest długi, ale w Prologu już nie, więc model poświęca więcej rozumowania mniej znanym
językom. Prolog jest najdroższym językiem: 26\,\% kosztu generowania przy jednej siódmej zadań.
Generowanie zajęło łącznie 42~minuty czasu pracy modelu (25,9\,s na wywołanie), a ewaluacja
w Dockerze dodaje 26--32\,s na przebieg dla Haskella, Scali i Prologu z powodu kompilacji
i uruchamiania.

\paragraph{Koszt} 98 wywołań generujących zużyło 501\,231 tokenów wejściowych i 635\,269 wyjściowych
i kosztowało \usd{0,368}. 70 wywołań sędziego zużyło 403\,703 tokeny wejściowe i 72\,482 wyjściowe
i kosztowało \usd{1,532}, około czterokrotnie więcej niż generowanie, ponieważ sędzia jest większym
modelem i czyta specyfikację, kontrakt oraz rubrykę. Cały przebieg kosztował około \usd{1,90}.
Powtórzone próbki są więc w zasięgu: dziesięć próbek na zadanie kosztowałoby około \usd{19}.

\paragraph{Wykonywanie instrukcji} Choć każdy prompt wymaga surowego kodu, 37 z 98 odpowiedzi
(37,8\,\%) otaczało kod znacznikami bloków Markdown: 28 z 70 pierwszych odpowiedzi (40\,\%) i 9 z 28
odpowiedzi naprawczych (32\,\%). Częstość ta silnie zależy od języka, od 0/13 dla Pascala i 1/15 dla
Haskella do 10/17 dla Prologu i 9/16 dla Javy. Nie wynika ze sformułowań szablonów. Szablony dla
Pascala i Prologu jako jedyne nie wymieniają wprost zakazanego znacznika (\code{no ```c blocks}),
a mimo to jeden z nich ma najniższą częstość, a drugi najwyższą. Znaczniki są nieszkodliwe, bo system
wycina kod. Pięć odpowiedzi z prozą bez znaczników (\cref{sec:failure-analysis}) już takie nie jest,
a metryka oparta na znacznikach ich nie liczy. Zgodność z formatem odpowiedzi najlepiej więc mierzyć
na dwa sposoby: przez znaczniki bloków i przez błędy kompilacji spowodowane tekstem niebędącym
kodem.
